\begin{abstract}
As semiconductor scaling approaches physical limits, the demand for unconventional computing architectures capable of solving intractable combinatorial optimization problems has intensified. This paper presents a rigorous theoretical framework and simulation validation for a Spatial Photonic Ising Machine (SPIM). By encoding spin variables into the phase of coherent light and leveraging the Fourier transform properties of free-space propagation, we realize a reconfigurable, fully-connected Ising model capable of finding ground states for NP-hard problems. We demonstrate that the system's scalability is governed by the bandwidth of Spatial Light Modulators (SLMs) rather than thermal dissipation, offering an $O(1)$ time-complexity advantage for calculating spin interactions in the optical domain. Furthermore, we analyze the industrial readiness of this technology, highlighting immediate applications in logistics, portfolio optimization, and computational biology, and argue that photonic computing represents a critical strategic pivot for next-generation high-performance computing infrastructure.
\end{abstract}
